

%\fancyhf{} % clear all header and footer fields
%\fancyhead[RO,R]{\thepage} %RO=right odd, RE=right even
%\renewcommand{\headrulewidth}{0pt}

\begin{center}
    %\Large
    \textbf{Argumentative Graph-RAG for Participatory Democracy}
    
    \vspace{0.4cm}
    
    \vspace{0.4cm}
    \textbf{Michele Pagano Mariano}
    
    \vspace{0.9cm}
    \textbf{Abstract}
    
\end{center}

Participatory democracy platforms generate massive volumes of textual data, leading to severe information overload for both citizens and decision-makers. While Natural Language Processing (NLP) and Argument Mining (AM) can extract argumentative structures, navigating this complex data remains a challenge.
    
This thesis explores the integration of Argument Mining with Retrieval-Augmented Generation (RAG) to create "Graph-RAG." By constructing knowledge graphs that represent the explicit relationships (support, attack, equivalence) between citizen proposals, we aim to enhance the ability of Large Language Models (LLMs) to retrieve context-rich information and generate balanced summaries.
    
We propose a pipeline that first builds argumentative graphs from civic datasets and then utilizes graph traversal to power a RAG system. This approach addresses the limitations of standard text-based retrieval by exploiting the structural logic of debates, ultimately providing intelligent tools for navigating large-scale deliberative discussions.